\documentclass[10pt,a4paper]{amsart}
\usepackage[margin=19mm]{geometry}
\usepackage{amsmath,amssymb,mathtools}
\usepackage[scr=rsfs,cal=euler]{mathalfa}
\usepackage{xfrac}
\usepackage[unicode,hidelinks,bookmarks=false]{hyperref}
\newcommand{\ie}{i.e.\ }
\newcommand{\cf}{cf.\ }
\newcommand{\resp}{respectively\ }
\newcommand{\Subsection}{\subsection}
\providecommand{\nobreakdash}{\nobreak\hbox{-}}
\setlength{\emergencystretch}{2em}
\multlinegap0pt
\usepackage{mathtools}
\newtheorem{theorem}{Theorem}
\newtheorem{lemma}[theorem]{Lemma}
\newtheorem{proposition}[theorem]{Proposition}
\newtheorem{corollary}[theorem]{Corollary}
\newtheorem{remark}[theorem]{Remark}
\numberwithin{equation}{section}
\DeclareMathOperator{\wt}{wt}
\DeclareMathOperator{\CT}{CT}
\DeclareMathOperator{\Newt}{Newt}
\DeclareMathOperator{\supp}{supp}
\newcommand{\C}{\mathbb{C}}
\newcommand{\Q}{\mathbb{Q}}
\newcommand{\Z}{\mathbb{Z}}
\newcommand{\N}{\mathbb{N}}
\newcommand{\E}{\mathbb{E}}
\newcommand{\Qbar}{\overline{\Q}}
\newcommand{\GMC}{\mathrm{GMC}}
\begin{document}
\title[A face-isolation proof of the two-variable Gaussian Moments ]{A face-isolation proof of the two-variable Gaussian Moments Conjecture}
\author{Michael Wilson}
\date{}
\maketitle
\noindent\emph{Source and licence.} A face-isolation proof of the two-variable Gaussian Moments Conjecture. Version arXiv:2607.23887v1 (CC BY 4.0). This material is distributed under the Creative Commons Attribution 4.0 International licence, http://creativecommons.org/licenses/by/4.0/, as recorded in the arXiv OAI licence metadata for this exact version and shown on the abstract page https://arxiv.org/abs/2607.23887v1. Full rights notice: licenses/arxiv-2607.23887-CC-BY-4.0.txt.\par\smallskip
\noindent\textit{Editorial scope.} Counted unit: the original \texttt{lemma} environment \texttt{lem:isolation} at source lines 264--275 together with its complete proof at lines 277--392; the delivered document keeps source lines 67--480 so that every referenced label and every citation resolves inside it. Native editable source is the paper's own concatenated TeX; the AMS wrapper is editorial formatting only. No publisher class, upstream script or unrelated artwork is redistributed.
\begin{equation}\label{eq:hyp}
  \E\bigl(P(X)^m\bigr)=0 \qquad (m\ge1),
\end{equation}
then for every $Q\in\C[x_1,\dots,x_n]$ one has $\E\bigl(Q(X)P(X)^m\bigr)=0$ for
all sufficiently large $m$. Equivalently, the space of polynomials with
vanishing Gaussian expectation is a Mathieu--Zhao space in the sense of Mathieu
\cite{Mathieu}. Complex coefficients are essential: for real $P$ the case $m=2$
already forces $P=0$. Derksen, van den Essen and Zhao proved that $\GMC(n)$ for
all $n$ implies the Jacobian Conjecture \cite[Thm.~1.6]{DvdEZ}.

The status of the conjecture changed in July 2026. It holds for $n=1$
\cite[Prop.~4.2]{DvdEZ}. In the opposite direction, Long \cite{Long} exhibited
explicit polynomials showing that $\GMC(n)$ fails for every $n\ge3$: a six-term
cubic in four variables settles $n\ge4$, and a five-term quartic in three
variables settles $n\ge3$. The only dimension left unresolved was $n=2$.

Throughout we work in the complex coordinates
\begin{equation}\label{eq:coords}
  Z=\frac{X+iY}{\sqrt2},\qquad W=\frac{X-iY}{\sqrt2},
\end{equation}
which form an invertible linear change of coordinates, so that
$\C[X,Y]=\C[Z,W]$ and total degree is the same in either. As random variables
$W=\overline{Z}$; both are complex linear polynomials in $X,Y$. We assign the
weight
\[
  \wt(Z^aW^b)=a-b,
\]
so that Gaussian expectation annihilates every monomial of nonzero weight; see
\eqref{eq:moments} below.

Two classes of $P$ were already known to satisfy the conclusion. If $P$ is
homogeneous, $\GMC(2)$ holds for $P$ \cite[Cor.~4.4]{DvdEZ}, by an argument
passing through the theorem of Duistermaat and van der Kallen on Laurent
polynomials \cite{DvdK}. If $P=ZA(U)+WB(U)$ with $U=ZW$, then Long
\cite[\S7]{Long} shows that \eqref{eq:hyp} forces $A=0$ or $B=0$. In both cases
the surviving polynomials are supported on monomials of a single sign of weight,
and Long concludes that a counterexample to $\GMC(2)$ would have to involve a
richer inhomogeneous mixture of positive and negative weights. The purpose of
this paper is to show that no such mixture exists.

\begin{theorem}\label{thm:main}
Let $X,Y$ be independent standard real Gaussian random variables and let
$P\in\C[X,Y]$ satisfy $\E(P^m)=0$ for every $m\ge1$. Write $P$ in the
coordinates \eqref{eq:coords}. Then either every monomial occurring in $P$ has
weight at least $1$, or every monomial occurring in $P$ has weight at most
$-1$. Consequently, for every $Q\in\C[X,Y]$,
\[
  \E\bigl(QP^m\bigr)=0 \qquad\text{whenever } m>\deg Q .
\]
In particular $\GMC(2)$ holds, with the explicit threshold $m\ge\deg Q+1$.
\end{theorem}

Together with $\GMC(1)$ \cite[Prop.~4.2]{DvdEZ} and \cite{Long}, this settles
the conjecture in every dimension: $\GMC(n)$ is true for $n\le2$ and false for
$n\ge3$.

The proof occupies Sections \ref{sec:prelim}--\ref{sec:proof}. Section
\ref{sec:prelim} records the moment evaluation and the elementary weight bound
that supplies the threshold. Section \ref{sec:spec} reduces to algebraic
coefficients; the point of that reduction, emphasised in Remark
\ref{rem:supportonly}, is that everything proved afterwards about the Newton
polygon is a statement about the \emph{support} of $P$ alone, so it transfers
back to the original complex polynomial. Section \ref{sec:isolation} contains
the main new ingredient, a prime isolation theorem for an arbitrary lower face
of the Newton polygon. Section \ref{sec:edge} shows that a Newton polygon
carrying weights of both signs always exposes a face to which prime isolation
applies, and Section \ref{sec:proof} assembles the proof. Section
\ref{sec:remarks} explains why the argument is specific to $n=2$ and records two
corollaries.

\section{Coordinates, moments and weights}\label{sec:prelim}

Throughout, $X$ and $Y$ are independent standard real Gaussians and $Z,W$ are as
in \eqref{eq:coords}. We use repeatedly the evaluation
\begin{equation}\label{eq:moments}
  \E\bigl(Z^aW^b\bigr)=\delta_{ab}\,a! \qquad (a,b\ge0).
\end{equation}
Indeed, the law of $(X,Y)$ is invariant under rotation, which multiplies
$Z^aW^b$ by $e^{i(a-b)\theta}$; this forces the expectation to vanish when
$a\ne b$. When $a=b$, the variable $ZW=(X^2+Y^2)/2$ is exponentially
distributed with mean one, whence $\E\bigl((ZW)^a\bigr)=a!$.

\begin{remark}
We stress that $Z$ and $W$ are \emph{not} an independent pair: as random
variables $W=\overline Z$, and $\E(Z^aW^a)=a!$ rather than
$\bigl((a-1)!!\bigr)^2$. Equation \eqref{eq:moments} is an identity about two
independent \emph{real} Gaussians written in complex linear coordinates; no
complex Gaussian measure is used anywhere in this paper.
\end{remark}

We write $\supp P\subset\N^2$ for the set of exponent pairs occurring in $P$,
and $\Newt(P)=\mathrm{conv}(\supp P)\subset\mathbb R_{\ge0}^2$ for its Newton
polygon.

\begin{lemma}[One-sided weights]\label{lem:weights}
Let $P\in\C[X,Y]$ be such that every monomial occurring in $P$ has weight at
least $1$, or such that every monomial occurring in $P$ has weight at most
$-1$. Then $\E(P^m)=0$ for every $m\ge1$, and for every $Q\in\C[X,Y]$,
\[
  \E\bigl(QP^m\bigr)=0 \qquad\text{whenever } m>\deg Q .
\]
\end{lemma}

\begin{proof}
Assume every monomial of $P$ has weight at least $1$. Weight is additive on
products, so every monomial occurring in $P^m$ has weight at least $m$. Write
$Q=\sum_{i,j}q_{ij}Z^iW^j$; each monomial of $Q$ has weight $i-j$ with
$|i-j|\le i+j\le\deg Q$. Hence every monomial of $QP^m$ has weight at least
$m-\deg Q$, which is strictly positive as soon as $m>\deg Q$. By
\eqref{eq:moments} Gaussian expectation annihilates each such monomial, so
$\E(QP^m)=0$. Taking $Q=1$ gives $\E(P^m)=0$ for $m\ge1$. The case of weights at
most $-1$ follows by exchanging $Z$ and $W$, which is the substitution
$Y\mapsto-Y$ and leaves the law of $(X,Y)$ invariant. The zero polynomial
satisfies both hypotheses vacuously.
\end{proof}

\section{Reduction to algebraic coefficients}\label{sec:spec}

\begin{lemma}[Algebraic specialization]\label{lem:spec}
Fix a finite set $S\subset\N^2$. Suppose there exists
$P=\sum_{\alpha\in S}c_\alpha Z^{a_\alpha}W^{b_\alpha}\in\C[Z,W]$ with
$c_\alpha\ne0$ for every $\alpha\in S$ and $\E(P^m)=0$ for every $m\ge1$. Then
there exists such a polynomial with all coefficients in $\Qbar$.
\end{lemma}

\begin{proof}
By \eqref{eq:moments}, for each $m\ge1$,
\begin{equation}\label{eq:momentexp}
  \E(P^m)=\sum_{\substack{k\in\N^{S},\ |k|=m\\ A(k)=B(k)}}
  \binom{m}{k}\Bigl(\prod_{\alpha\in S}c_\alpha^{k_\alpha}\Bigr)A(k)!,
  \qquad
  A(k)=\sum_\alpha k_\alpha a_\alpha,\quad B(k)=\sum_\alpha k_\alpha b_\alpha,
\end{equation}
which is a polynomial $\varphi_m\in\Z[c_\alpha:\alpha\in S]$ in the
coefficients. Let $I\subset\Q[c_\alpha]$ be the ideal generated by
$\{\varphi_m\}_{m\ge1}$; it is finitely generated by the Hilbert basis theorem.
Let $f=\prod_{\alpha\in S}c_\alpha$. The set
$V(I)\cap D(f)$ is a constructible set defined over $\Q$, and by hypothesis it
has a $\C$-point. Since $\Qbar$ is algebraically closed and
$\Qbar\subset\C$, the Nullstellensatz gives a $\Qbar$-point of the same set,
which is the required polynomial.
\end{proof}

\begin{remark}[Why this reduction is legitimate]\label{rem:supportonly}
Lemma \ref{lem:spec} produces a \emph{different} polynomial, so it is worth
recording precisely how it will be used. Step 1 of Section \ref{sec:proof}
argues by contradiction: from a complex $P$ with support $S$ carrying weights of
both signs and satisfying \eqref{eq:hyp}, Lemma \ref{lem:spec} supplies
$P'\in\Qbar[Z,W]$ with the same support and the same vanishing, and Lemmas
\ref{lem:isolation} and \ref{lem:edge} applied to $P'$ then contradict each
other. Both of those conclusions are assertions about the exposed face
$F\subset S$ alone ,  which weights occur on $F$ ,  and mention no
coefficients, so the contradiction is not attached to the particular
$\Qbar$-point produced. What transfers the conclusion back to $P$ is simply the
contrapositive of Lemma \ref{lem:spec}: no polynomial over $\Qbar$ with support
$S$ satisfies \eqref{eq:hyp}, hence none over $\C$ does either. This is the only
use we make of Lemma \ref{lem:spec}.
\end{remark}

\section{Prime isolation of a lower face}\label{sec:isolation}

Let $P\in\Qbar[Z,W]$ with $\supp P=S$, let $K=\Q(c_\alpha:\alpha\in S)$, a
number field, and let $\mathcal O_K$ be its ring of integers. For an integral
linear form
\[
  \ell(a,b)=ca+db, \qquad c,d\in\Z,\quad c+d>0,
\]
put $\lambda=\min_{\alpha\in S}\ell(\alpha)$ and let
$F=\{\alpha\in S:\ell(\alpha)=\lambda\}$ be the exposed face. We may and do
assume $\gcd(c,d)=1$: replacing $(c,d)$ by $(c,d)/\gcd(c,d)$ scales $\lambda$
and $c+d$ by the same factor, leaves $F$ and $\lambda/(c+d)$ unchanged, and only
shrinks the set of primes excluded below. Define the \emph{face weight
polynomial}
\[
  R(z)=\sum_{\alpha\in F}c_\alpha z^{a_\alpha-b_\alpha}
  \ \in\ K[z,z^{-1}].
\]

\begin{remark}[No cancellation in $R$]\label{rem:injective}
The weight $w(a,b)=a-b$ is injective on $F$. Indeed, if $\alpha,\alpha'\in F$
satisfy $w(\alpha)=w(\alpha')$, then $v=\alpha-\alpha'$ lies in
$\ker\ell\cap\ker w$, so $cv_1+dv_2=0$ and $v_1=v_2$, whence $(c+d)v_1=0$ and
$v=0$ because $c+d>0$. Equivalently, $\ker\ell=\langle(d,-c)\rangle$ and
$\ker w=\langle(1,1)\rangle$ meet only in the origin, since
$\det\begin{psmallmatrix}d&-c\\1&1\end{psmallmatrix}=c+d\ne0$.

Consequently $R$ has exactly $|F|$ terms, with pairwise distinct exponents and
with the nonzero coefficients $c_\alpha$: no cancellation can occur, and the
map $\alpha\mapsto w(\alpha)$ is a bijection from $F$ onto the exponent set of
$R$. This is what licenses the passage, in Lemma \ref{lem:isolation}, from a
statement about the exponents of $R$ to a statement about the weights occurring
on $F$. Without it a weight-zero point of $F$ could in principle be cancelled
out of $R$ by another face point of the same weight, and the contradiction in
Step 1 of Section \ref{sec:proof} would fail in exactly the sub-case where $w$
vanishes somewhere on $F$.
\end{remark}


\begin{lemma}[Prime isolation]\label{lem:isolation}
With the notation above, suppose $\E(P^m)=0$ for every $m\ge1$ and
$\lambda/(c+d)\ge0$. Let $B\ge1$ be the denominator of $\lambda/(c+d)$ in lowest
terms. Then
\begin{equation}\label{eq:CTvanish}
  \CT R(z)^{Bn}=0 \qquad (n\ge1).
\end{equation}
Consequently every exponent occurring in $R$ is strictly positive, or every
exponent occurring in $R$ is strictly negative. By Remark \ref{rem:injective}
this says equivalently that all weights occurring on $F$ have the same strict
sign.
\end{lemma}

\begin{proof}
Fix $n\ge1$ and set $q=Bn$, so that
\[
  R_0:=\frac{q\lambda}{c+d}\in\N .
\]
Call a rational prime $p$ \emph{good} (for the data $P,\ell,q$) if
\begin{enumerate}
\item[(G1)] $p$ is unramified in $K$;
\item[(G2)] every $c_\alpha$ is $\mathfrak p$-integral for every prime
$\mathfrak p$ of $\mathcal O_K$ above $p$;
\item[(G3)] $p\nmid(c+d)$;
\item[(G4)] $p>R_0$.
\end{enumerate}
All but finitely many primes satisfy (G1)--(G3), and (G4) excludes finitely many
more once $q$ is fixed; so infinitely many good primes exist. Note that (G2) asks
only for integrality, not for the reductions $\overline{c_\alpha}$ to be
nonzero: the moment computation below takes place in characteristic zero and is
reduced only at the very end, so a coefficient meeting $\mathfrak p$ does no
harm. Fix a good $p$, fix $\mathfrak p\mid p$, and let $v=v_{\mathfrak p}$ be
normalised by $v(p)=1$, which is possible by (G1). Set $m=qp$.

Consider the expansion \eqref{eq:momentexp} of $\E(P^m)$. For a multi-index $k$
occurring there we have $A(k)=B(k)=:N$, and therefore
\begin{equation}\label{eq:balance}
  (c+d)N=cA(k)+dB(k)=\sum_{\alpha}k_\alpha\ell(\alpha).
\end{equation}
Since $\ell(\alpha)\ge\lambda$ for every $\alpha\in S$, \eqref{eq:balance} gives
$(c+d)N\ge m\lambda=qp\lambda$, hence
\begin{equation}\label{eq:Nlower}
  N\ \ge\ pR_0 .
\end{equation}

\smallskip
\noindent\emph{The pure-face contribution.} Suppose $k_\alpha=0$ for every
$\alpha\notin F$. Then \eqref{eq:balance} reads $(c+d)N=qp\lambda$, so
$N=pR_0$, independently of $k$. Moreover the balance condition $A(k)=B(k)$ is
exactly $\sum_\alpha k_\alpha(a_\alpha-b_\alpha)=0$, which is the condition
selecting the constant term of $R^{m}$. Hence the sum of all pure-face terms in
\eqref{eq:momentexp} equals
\begin{equation}\label{eq:pureface}
  (pR_0)!\,\CT R(z)^{qp}.
\end{equation}

\smallskip
\noindent\emph{Every other contribution is divisible by $p^{R_0+1}$.} Suppose
$k_\beta>0$ for some $\beta\notin F$. We distinguish two cases.

\emph{Case 1: $p\nmid k_\alpha$ for some $\alpha$.} The base-$p$ units digit of
$m=qp$ is $0$, while the units digits of the $k_\alpha$ sum to a multiple of
$p$ which is positive, hence at least $p$. Adding the $k_\alpha$ in base $p$
therefore produces a carry, and by Kummer's theorem for multinomial
coefficients, $v\bigl(\binom{m}{k}\bigr)\ge1$. By \eqref{eq:Nlower},
\[
  v(N!)\ \ge\ \Bigl\lfloor\frac{N}{p}\Bigr\rfloor\ \ge\ R_0 ,
\]
so the term has valuation at least $R_0+1$.

\emph{Case 2: $k_\alpha=p\,l_\alpha$ for every $\alpha$.} Then
$\sum_\alpha l_\alpha=q$, and
\[
  \Delta:=\sum_\alpha l_\alpha\bigl(\ell(\alpha)-\lambda\bigr)
\]
is a positive integer, since $\ell(\beta)>\lambda$ and $l_\beta>0$. Now
\eqref{eq:balance} gives $(c+d)N=p(q\lambda+\Delta)$. As $q\lambda=R_0(c+d)$ is
divisible by $c+d$, we get $(c+d)\mid p\Delta$, and $p\nmid(c+d)$ by (G3),
whence $(c+d)\mid\Delta$ and $\Delta/(c+d)\ge1$. Therefore
\[
  N=p\Bigl(R_0+\frac{\Delta}{c+d}\Bigr)\ \ge\ p(R_0+1),
  \qquad
  v(N!)\ \ge\ \Bigl\lfloor\frac{N}{p}\Bigr\rfloor\ \ge\ R_0+1 ,
\]
so again the term has valuation at least $R_0+1$.

\smallskip
\noindent\emph{Conclusion of the valuation count.} By hypothesis
$\E(P^{qp})=0$. Subtracting the contributions just bounded, \eqref{eq:pureface}
gives
\[
  (pR_0)!\,\CT R^{qp}\ \equiv\ 0 \pmod{\mathfrak p^{R_0+1}} .
\]
Since $p>R_0$ by (G4), we have
$v\bigl((pR_0)!\bigr)=\lfloor pR_0/p\rfloor=R_0$, and dividing by $p^{R_0}$
yields
\[
  \CT R^{qp}\ \equiv\ 0 \pmod{\mathfrak p} .
\]
Let $\kappa=\mathcal O_K/\mathfrak p$, a field of characteristic $p$, and let
$\overline R\in\kappa[z,z^{-1}]$ be the reduction of $R$, which is defined by
(G2). Writing $\overline R=\sum_i \bar a_iz^i$ we have
$\overline R^{\,p}=\sum_i\bar a_i^{\,p}z^{ip}$, whence
\[
  \CT \overline R^{\,qp}
  =\CT\bigl(\overline R^{\,q}\bigr)^{p}
  =\bigl(\CT \overline R^{\,q}\bigr)^{p} .
\]
(We emphasise that the correct identity is $\CT(g^p)=(\CT g)^p$; the variant
$\CT(g^p)=\CT(g)$ is valid only when the coefficients lie in the prime field,
which need not hold here.) Since $\kappa$ is a field, it follows that
$\CT R^{q}\equiv0\pmod{\mathfrak p}$. This holds for infinitely many good
primes $p$, and an element of $K$ lying in infinitely many distinct primes of
$\mathcal O_K$ is zero; hence $\CT R^{q}=0$. As $n\ge1$ was arbitrary and
$q=Bn$, this proves \eqref{eq:CTvanish}.

\smallskip
\noindent\emph{Application of Duistermaat--van der Kallen.} Put $g=R^{B}$, a
Laurent polynomial in one variable. By \eqref{eq:CTvanish},
$\CT(g^{n})=\CT R^{Bn}=0$ for every $n\ge1$. By the theorem of Duistermaat and
van der Kallen \cite{DvdK}, a Laurent polynomial in one variable which is
neither a polynomial in $z$ nor a polynomial in $z^{-1}$ has some positive power
with nonzero constant term; combined with $\CT(g)=0$ this forces
$g\in z\C[z]$ or $g\in z^{-1}\C[z^{-1}]$, equivalently $0\notin\Newt(g)$. Since
the extreme coefficients of $R^{B}$ are powers of the extreme coefficients of
$R$ and hence nonzero, $\Newt(R^{B})=B\cdot\Newt(R)$, and as $B\ge1$ we conclude
$0\notin\Newt(R)$. In one variable this says exactly that every exponent of $R$
is strictly positive or every exponent is strictly negative.
\end{proof}


\section{The mixed-edge lemma}\label{sec:edge}

\begin{lemma}[Mixed edge]\label{lem:edge}
Let $S\subset\N^2$ be finite and $\Delta=\mathrm{conv}(S)\subset\mathbb
R_{\ge0}^2$. Suppose the weight $w(a,b)=a-b$ takes both a positive and a
negative value on $S$. Then there is an integral linear form
$\ell(a,b)=ca+db$ with $c+d>0$ such that, writing
$\lambda=\min_{\alpha\in S}\ell(\alpha)$ and $F=\{\alpha\in S:\ell(\alpha)=\lambda\}$,
\begin{enumerate}
\item[(i)] $\lambda/(c+d)\ge0$, and
\item[(ii)] $w$ does not have one strict sign on $F$: either $w$ vanishes
somewhere on $F$, or $w$ takes both signs on $F$.
\end{enumerate}
\end{lemma}

\begin{proof}
\emph{Case 1: $\Delta$ is one-dimensional}, i.e.\ $S$ is contained in a line.
Let $(\Delta a,\Delta b)$ be a direction vector of that line, taken to be an
integral vector, and set $(c,d)=\pm(-\Delta b,\Delta a)$ with the sign chosen so
that $c+d=\pm(\Delta a-\Delta b)>0$; this is possible because $w$ is not
constant on $S$, so $\Delta a-\Delta b\ne0$. Then $\ell$ is constant on the
line, so $F=S$ and (ii) holds by hypothesis. For (i): $w$ takes both signs on
the segment $\Delta$ and is affine there, so $\Delta$ meets the diagonal at some
point $(t,t)$; since $\Delta\subset\mathbb R_{\ge0}^2$ we have $t\ge0$, and
$\lambda=\ell(t,t)=(c+d)t$, whence $\lambda/(c+d)=t\ge0$.

\emph{Case 2: $\Delta$ is two-dimensional.} The maximum of $w$ over $\Delta$
equals its maximum over $S$, which is positive, and is attained at a vertex;
similarly some vertex has $w<0$. Traversing the boundary of $\Delta$
counterclockwise, the continuous piecewise affine function $w$ therefore takes
both signs, so there is an edge $E$ from $p$ to $q$ (in the counterclockwise
order) with $w(p)\le0\le w(q)$ and $w(p)<w(q)$. Let
$(\Delta a,\Delta b)=q-p$, an integral vector, and put
$(c,d)=(-\Delta b,\Delta a)$, dividing by $\gcd$ to make it primitive. For a
counterclockwise-oriented convex polygon, $(-\Delta b,\Delta a)$ is the inward
normal of $E$, so $\ell$ attains its minimum over $\Delta$ exactly on $E$; in
particular $F=S\cap E$ and $\lambda=\ell(E)$. Its coordinate sum is
$c+d=\Delta a-\Delta b=w(q)-w(p)>0$. Since $w$ is affine on $E$ with
$w(p)\le0\le w(q)$, the edge meets the diagonal at some $(t,t)\in E$, and
$t\ge0$ because $E\subset\mathbb R_{\ge0}^2$; hence
$\lambda=\ell(t,t)=(c+d)t$ and $\lambda/(c+d)=t\ge0$, which is (i). Finally the
endpoints $p,q$ are vertices of $\Delta$, hence lie in $S$, hence lie in $F$;
as $w(p)\le0\le w(q)$, statement (ii) holds.
\end{proof}

\section{Proof of Theorem \ref{thm:main}}\label{sec:proof}

Let $P\in\C[X,Y]$ satisfy $\E(P^m)=0$ for every $m\ge1$, and let $S=\supp P$ in
the coordinates \eqref{eq:coords}. If $P=0$ there is nothing to prove, so assume
$S\ne\emptyset$.

\smallskip
\noindent\emph{Step 1: $S$ does not carry weights of both signs.} Suppose it
did. By Lemma \ref{lem:spec} there is $P'\in\Qbar[Z,W]$ with $\supp P'=S$ and
$\E(P'^m)=0$ for every $m\ge1$. Apply Lemma \ref{lem:edge} to $S$, obtaining an
integral form $\ell$ with $c+d>0$ and $\lambda/(c+d)\ge0$ whose exposed face $F$
carries a weight equal to zero, or weights of both signs. Apply Lemma
\ref{lem:isolation} to $P'$ and this $\ell$: all weights occurring on $F$ have
the same strict sign, the passage from the exponents of $R$ to the weights on
$F$ being justified by Remark \ref{rem:injective}. These two assertions about
$F$ contradict each other. Hence no polynomial over $\Qbar$ with support $S$
satisfies \eqref{eq:hyp}, and by Lemma \ref{lem:spec} none over $\C$ does
either (Remark \ref{rem:supportonly}). So every weight occurring in $P$ is
nonnegative, or every weight occurring in $P$ is nonpositive.

\smallskip
\noindent\emph{Step 2: the weight-zero part vanishes.} A monomial of weight
zero is $Z^aW^a=U^a$ with $U=ZW$, so the weight-zero part of $P$ is $C(U)$ for
some $C\in\C[U]$. By Step 1 all other monomials of $P$ have weight of one strict
sign, so a product of $m$ monomials of $P$ has weight zero only if every factor
is drawn from $C(U)$; hence the weight-zero part of $P^m$ is exactly $C(U)^m$.
By \eqref{eq:moments},
\[
  0=\E(P^m)=\mathcal L(C^m),\qquad \mathcal L(U^j)=j! \qquad (m\ge1).
\]
The one-variable case of the Factorial Conjecture, proved by van den Essen,
Wright and Zhao \cite[Thm.~4.9]{vdEWZ}, states that $\mathcal L(C^m)=0$ for
every $m\ge1$ forces $C=0$. Hence $P$ has no weight-zero monomial, and by Step
1 every monomial of $P$ has weight at least $1$, or every monomial has weight at
most $-1$.

\smallskip
\noindent\emph{Step 3: mixed moments.} Lemma \ref{lem:weights} now gives
$\E(QP^m)=0$ for every $Q\in\C[X,Y]$ and every $m>\deg Q$. \qed

\section{Remarks and corollaries}\label{sec:remarks}


\begin{thebibliography}{99}
\bibitem[DvdEZ]{DvdEZ} H.~Derksen, A.~van den Essen and W.~Zhao, \emph{The Gaussian Moments Conjecture and the Jacobian Conjecture}, Israel J. Math. \textbf{219} (2017), no.~2, 917--928. \href{https://arxiv.org/abs/1506.05192}{arXiv:1506.05192}.
\bibitem[DvdK]{DvdK} J.~J.~Duistermaat and W.~van der Kallen, \emph{Constant terms in powers of a Laurent polynomial}, Indag. Math. (N.S.) \textbf{9} (1998), no.~2, 221--231.
\bibitem[Long]{Long} C.~D.~Long, \emph{Small counterexamples to the Gaussian Moments Conjecture}, preprint (2026). \href{https://arxiv.org/abs/2607.18186}{arXiv:2607.18186}.
\bibitem[Mathieu]{Mathieu} O.~Mathieu, \emph{Some conjectures about invariant theory and their applications}, in: Alg\`ebre non commutative, groupes quantiques et invariants (Reims, 1995), S\'emin. Congr. \textbf{2}, Soc. Math. France, Paris, 1997, pp.~263--279.
\bibitem[vdEWZ]{vdEWZ} A.~van den Essen, D.~Wright and W.~Zhao, \emph{On the Image Conjecture}, J. Algebra \textbf{340} (2011), 211--224. \href{https://arxiv.org/abs/1008.3962}{arXiv:1008.3962}.
\end{thebibliography}
\end{document}
